\documentclass[preprint,12pt]{elsarticle}

\usepackage{amsmath,amssymb}
\usepackage{booktabs}
\usepackage{graphicx}
\usepackage{xcolor}
\usepackage{hyperref}

\journal{Computer-Aided Design}

\newcommand{\TODO}[1]{\textcolor{red}{[TODO: #1]}}
\newcommand{\R}{\mathbb{R}}
\newcommand{\Kmax}{K_{\max}}

\begin{document}

\begin{frontmatter}

\title{Canonical Ordinal Model-Order Prediction for Learning-Based
B-Spline Curve Fitting}

% Computer-Aided Design uses double-anonymized review. Author information is
% intentionally kept in title_page.tex rather than this manuscript.

\begin{abstract}
Recovering a compact parametric B-spline from ordered point samples requires
estimating both a discrete model order and continuous knot locations. Existing
learned pruning formulations often use independent or weakly coupled binary
gates, for which a fixed threshold must simultaneously determine how many knots
to retain and which knots are geometrically necessary. This coupling can yield
poorly calibrated structures despite low curve-fitting error. We formulate the
problem instead as tolerance-aware canonical model-order selection followed by
count-conditioned knot prediction. Redundant source knots are removed under a
fixed geometric tolerance to construct identifiable supervision. A local
cross-attention head predicts the internal-knot count through ordinal targets,
and a shared count-conditioned query decoder returns exactly that many ordered
knot locations. Complete count branches are compared using a learned-prior BIC
criterion. The resulting knot vector is converted
to a standard open cubic B-spline whose control points are obtained by regularized
least squares. \TODO{Insert final quantitative improvements on count accuracy,
knot F1, matched knot error, fitting RMS, and runtime.} Experiments on synthetic
nonuniform cubic B-splines and \TODO{a real or public CAD benchmark} evaluate
structural accuracy, geometric fidelity, robustness, and generalization. The
proposed decomposition removes deployment-time threshold tuning and provides a
direct, interpretable connection between learned inference and CAD model order.
\end{abstract}

\begin{keyword}
B-spline \sep knot prediction \sep curve fitting \sep geometric deep learning
\sep model order selection \sep computer-aided design
\end{keyword}

\end{frontmatter}

\section{Introduction}
\label{sec:introduction}

B-spline curves are a basic representation in geometric modeling because they
combine local control, compact storage, stable evaluation, and compatibility
with established CAD kernels \cite{deboor1972,cox1972,piegl1997}. Converting
ordered measurements into a compact B-spline nevertheless remains difficult:
point parameters, the number and locations of internal knots, and control points
are mutually dependent. Excess knots improve data fidelity but reduce
compactness and may fit noise, whereas too few knots erase geometric detail.

Learning-based inference offers a way to amortize this nonlinear estimation
across a family of curves. A common design predicts a fixed set of candidate
knots and attaches a binary existence gate to every candidate. Our preliminary
experiments show a structural failure of this strategy: a model can achieve a
small fitting error while retaining all candidates for many curves and none for
others. This occurs because a fixed probability threshold is asked to infer
both model order and candidate identity.

We instead define a parsimonious target by removing source knots under a fixed
geometric tolerance, treat the resulting number of internal knots as an ordinal
variable, and predict knot locations conditional on that number. This view
matches the semantics of CAD model selection: first choose the representation
complexity, then estimate its continuous parameters.

The intended contributions are:
\begin{itemize}
  \item a tolerance-aware canonical labeling procedure that avoids supervising
  arbitrary redundant source representations;
  \item a local-attention ordinal count estimator that removes binary candidate
  thresholding from B-spline knot inference;
  \item a shared count-conditioned query decoder that predicts exactly the
  selected number of ordered internal knots;
  \item an end-to-end fitting pipeline that converts the prediction to a
  standard open cubic B-spline through a regularized linear solve and BIC model
  selection; and
  \item a structural evaluation protocol that separates knot-count accuracy,
  knot localization, and geometric fitting fidelity.
\end{itemize}

\TODO{Rewrite the contribution list after the final implementation and ensure
that every claimed contribution is supported by an experiment.}

\section{Related work}
\label{sec:related_work}

\subsection{B-spline parameterization and knot placement}

Classical B-spline fitting alternates or jointly optimizes point parameters,
knot vectors, and control points. Standard references describe basis evaluation,
least-squares fitting, and practical knot construction
\cite{deboor1972,cox1972,piegl1997,dierckx1993}. Parameterization rules such as
chord length and centripetal parameterization remain strong deterministic
baselines \cite{lee1989}. \TODO{Add recent adaptive knot-placement and model
selection papers from Computer-Aided Design and Computer Aided Geometric Design.}

\subsection{Learning geometric representations for CAD}

Recent geometric learning methods infer parametric primitives, CAD structure,
or editable representations from sampled geometry. Large CAD datasets have
enabled supervised learning of geometric quantities and reconstruction
pipelines \cite{koch2019abc}. \TODO{Add the most relevant curve- and
spline-reconstruction methods, clearly distinguishing unordered point-cloud
methods from the ordered-curve setting studied here.}

\subsection{Sparse gates and set prediction}

Hard-Concrete gates provide a differentiable approximation to an $L_0$
penalty \cite{louizos2018}. Set-prediction architectures instead associate
learned queries with targets through assignment losses \cite{carion2020}.
Although these mechanisms are useful for variable-cardinality prediction,
independent gates require probability calibration and a deployment threshold.
Our formulation makes cardinality explicit and conditions continuous knot
prediction on the selected count.

\section{Problem formulation}
\label{sec:problem}

Let $Q=(q_i)_{i=0}^{M-1}\in\R^{M\times d}$ be ordered samples from an unknown
curve. We seek a degree-$p$ open B-spline
\begin{equation}
  C(t)=\sum_{r=0}^{n-1}N_{r,p}(t;U)P_r,
  \qquad t\in[0,1],
\end{equation}
with control points $P_r\in\R^d$ and knot vector
\begin{equation}
  U=(\underbrace{0,\ldots,0}_{p+1},
  u_1,\ldots,u_K,
  \underbrace{1,\ldots,1}_{p+1}),
  \quad 0<u_1<\cdots<u_K<1.
\end{equation}
The unknowns comprise point parameters
$0=t_0<\cdots<t_{M-1}=1$, the discrete internal-knot count
$K\in\{0,\ldots,\Kmax\}$, the knot locations, and the control points.

Given predictions $(\widehat t,\widehat K,\widehat U)$, the control points are
obtained by
\begin{equation}
  \widehat P=\arg\min_P
  \|A(\widehat t,\widehat U)P-Q\|_F^2
  +\lambda_s\|D_2P\|_F^2+\lambda_r\|P\|_F^2,
  \label{eq:control_solve}
\end{equation}
where $A$ is the B-spline collocation matrix and $D_2$ penalizes second
differences of the control polygon. The primary goal is not only low fitting
error but recovery of the simplest structurally correct representation.

\section{Method}
\label{sec:method}

\subsection{Geometry encoding and point parameterization}

The encoder receives point coordinates together with first and second
derivatives normalized by chord-length parameter gaps. It produces local tokens $F\in\R^{M\times h}$ and a pooled curve
descriptor $g\in\R^h$. A parameter head predicts positive intervals, which are
normalized to form a strictly increasing parameter sequence. A positional
encoding of the predicted parameters is added to the local tokens before knot
decoding.

\subsection{Canonical structural supervision}

Starting from a generated source knot vector, we greedily remove the single
knot whose removal gives the smallest refitting error. A removal is accepted
when the normalized Euclidean RMS remains below a prescribed tolerance
$\epsilon$. The terminal representation defines the canonical count $K^*$ and
knot locations $U^*$. Thus the target describes required geometric complexity
rather than the arbitrary representation used to generate the samples.

\subsection{Ordinal knot-count prediction}

A set of count queries cross-attends to the position-encoded local tokens and
produces a scalar complexity evidence $e$. Ordered thresholds $b_r$ define
\begin{equation}
  s_r=P(K\ge r)=\sigma(e-b_r),\qquad r=1,\ldots,\Kmax.
\end{equation}
Adjacent survival probabilities define a normalized class distribution. The
head is trained using binary ordinal targets $\mathbf 1[K^*\ge r]$ and an
asymmetric expected-overcount penalty.

\subsection{Count-conditioned knot decoder}

All admissible counts share interval queries, cross-attention, and output
weights. A learned count embedding conditions the first $K+1$ interval queries.
Positive interval logits are normalized to obtain exactly $K$ strictly ordered
internal knots. During training, the representation indexed by the true count
receives the knot-position and fitting losses.
The direct neural prediction used for validation is
\begin{equation}
  \widehat K=\arg\max_K p_K,
  \qquad \widehat U=\mathcal D_{\widehat K}(F,g).
\end{equation}
The deployment refinement described below may instead select among complete
count branches. No candidate deletion, binary gate, or probability threshold
is used.

\subsection{Training objective}

For a sample with count $K^*$ and true internal knots $U^*$, the objective is
\begin{equation}
  \mathcal L=\lambda_f\mathcal L_{fit}
  +\lambda_t\mathcal L_t
  +\lambda_K\mathcal L_{count}
  +\lambda_o\mathcal L_{over}
  +\lambda_u\mathcal L_{knot},
\end{equation}
where
\begin{align}
  \mathcal L_{fit}&=M^{-1}\sum_i\|C(\widehat t_i)-q_i\|_2^2,\\
  \mathcal L_t&=M^{-1}\sum_i(\widehat t_i-t_i^*)^2,\\
  \mathcal L_{count}&=\Kmax^{-1}\sum_r
  \operatorname{BCEWithLogits}(e-b_r,\mathbf 1[K^*\ge r]),\\
  \mathcal L_{over}&=\max(0,\mathbb E[K]-K^*),\\
  \mathcal L_{knot}&=(K^*)^{-1}\sum_{j=1}^{K^*}
  \operatorname{SmoothL1}(\widehat u_j^{(K^*)}-u_j^*).
\end{align}

\subsection{Standard B-spline deployment}

Every complete count-conditioned representation defines an open knot vector.
Equation \eqref{eq:control_solve} is solved for each admissible count, after
which BIC augmented by the learned count prior selects the deployed model.
This compares complete spline orders and does not delete individual candidates.

\begin{figure}[t]
  \centering
  \fbox{\parbox[c][35mm][c]{0.90\linewidth}{\centering
  \TODO{Architecture figure: ordered samples $\rightarrow$ encoder and
  parameter head $\rightarrow$ CountHead $\rightarrow$ selected
  count-conditioned decoder $\rightarrow$ standard B-spline refit.}}}
  \caption{Overview of the proposed count-conditioned fitting pipeline.}
  \label{fig:architecture}
\end{figure}

\section{Experimental setup}
\label{sec:experiments}

\subsection{Datasets}

The synthetic benchmark contains open cubic B-splines generated from smooth
random control polygons. The current configuration uses 64 ordered samples per
curve, 5--10 source control points, 1--6 source internal knots, nonuniform knot intervals,
nonuniform parameter sampling, and Gaussian coordinate noise with standard
deviation $10^{-3}$. Curves are centered and scaled to unit radius. Canonical
targets are produced by greedy knot removal under normalized RMS tolerance
$5\times10^{-3}$ and contain 0--6 internal knots. The default training and validation sets contain 10,000 and
1,000 curves with disjoint random seeds.

\TODO{Add an unseen test set and a public or real CAD dataset. Explain how
ground-truth curve parameterizations are made comparable.}

\subsection{Baselines and ablations}

The comparison should include:
\begin{enumerate}
  \item fixed-$\Kmax$ fitting without deletion;
  \item independent supervised existence gates;
  \item candidate self-attention with Hard-Concrete gates (v3);
  \item a CountHead followed by top-$K$ candidate ranking; and
  \item categorical count-conditioned decoding with source labels (v4);
  \item ordinal shared count-conditioned decoding with canonical labels; and
  \item the full method with BIC and learned count prior.
\end{enumerate}
Classical chord-length fitting and an optimization-based adaptive knot method
should be included when implementations are available. Ablations should remove
true-parameter supervision, position encoding, local cross-attention, count
conditioning, canonical labeling, decoder sharing, BIC selection, and
control-polygon regularization one at a time.

\subsection{Metrics}

We report count accuracy, count mean absolute error, knot precision/recall/F1
within tolerances $\{0.01,0.02,0.05\}$, matched-knot MAE, pointwise Euclidean
RMS, maximum deviation, control-point count, zero/all-knot failure rates, and
runtime. Knot matching is reported only under a shared parameterization. Model
selection uses the validation set; the test set is evaluated once.

\subsection{Implementation details}

\TODO{Record hardware, software versions, parameter count, optimizer, learning
rates, batch size, epoch count, random seeds, loss weights, and training time.
Report mean and standard deviation over at least three seeds.}

\section{Results and discussion}
\label{sec:results}

\subsection{Why threshold pruning fails}

Table~\ref{tab:v3diagnostic} records the observed v3 validation behavior. The
mean retained count is only moderately above the truth, but its distribution is
pathological: 49 of 128 curves retain all eight candidates and 22 retain none.
The result confirms that a mean-count loss does not constrain per-curve
cardinality and that globally coupled gate scores are poorly calibrated for a
fixed threshold.

\begin{table}[t]
  \centering
  \caption{Diagnostic results for the v3 threshold-pruning baseline on 128
  validation curves. These are not results of the proposed method.}
  \label{tab:v3diagnostic}
  \begin{tabular}{lr}
    \toprule
    Metric & Value \\
    \midrule
    True internal-knot count, mean & 3.281 \\
    Retained count at threshold 0.5, mean & 4.242 \\
    Curves retaining all 8 candidates & 49/128 \\
    Curves retaining zero candidates & 22/128 \\
    Knot precision / recall / F1 at 0.05 & 0.407 / 0.526 / 0.459 \\
    Matched-knot MAE & 0.0240 \\
    Standard B-spline refit RMS & 0.0160 \\
    \bottomrule
  \end{tabular}
\end{table}

\subsection{Main comparison}

The v4 categorical model reached a knot precision of approximately 0.44, but
its training and validation count accuracies were approximately 0.64 and 0.21,
respectively. Supplying the true count increased knot precision to approximately
0.55, whereas additionally supplying true point parameters produced no further
improvement. These diagnostics motivate canonical labels, ordinal local count
attention, and parameter sharing across count-conditioned representations.

\begin{table*}[t]
  \centering
  \caption{Main test-set comparison. Best values should be bold and the second
  best underlined only after all runs are complete.}
  \label{tab:main}
  \begin{tabular}{lrrrrrr}
    \toprule
    Method & Count acc. $\uparrow$ & Count MAE $\downarrow$ & Knot F1 $\uparrow$
    & Knot MAE $\downarrow$ & RMS $\downarrow$ & Time (ms) $\downarrow$ \\
    \midrule
    Fixed-$\Kmax$ & -- & \TODO{x} & \TODO{x} & \TODO{x} & \TODO{x} & \TODO{x} \\
    Independent gates & \TODO{x} & \TODO{x} & \TODO{x} & \TODO{x} & \TODO{x} & \TODO{x} \\
    Hard-Concrete v3 & \TODO{x} & \TODO{x} & \TODO{x} & \TODO{x} & \TODO{x} & \TODO{x} \\
    Count + top-$K$ & \TODO{x} & \TODO{x} & \TODO{x} & \TODO{x} & \TODO{x} & \TODO{x} \\
    Categorical count v4 & \TODO{x} & \TODO{x} & \TODO{x} & \TODO{x} & \TODO{x} & \TODO{x} \\
    Proposed ordinal, argmax & \TODO{x} & \TODO{x} & \TODO{x} & \TODO{x} & \TODO{x} & \TODO{x} \\
    Proposed ordinal, BIC & \TODO{x} & \TODO{x} & \TODO{x} & \TODO{x} & \TODO{x} & \TODO{x} \\
    \bottomrule
  \end{tabular}
\end{table*}

\TODO{Discuss whether improvements come from count accuracy, better knot
localization, or both. Include confidence intervals and statistical tests where
appropriate. Do not infer structural superiority from fitting error alone.}

\subsection{Robustness, generalization, and failure cases}

Evaluate increasing noise, nonuniform sampling, knot clustering, unseen control
point counts, and 3D curves. Show representative successes and failures,
including curves whose count is misclassified by one or more knots. Discuss
rank deficiency and the effect of control-polygon regularization.

\begin{figure}[t]
  \centering
  \fbox{\parbox[c][42mm][c]{0.90\linewidth}{\centering
  \TODO{Qualitative comparison: samples, true curve/knots, v3 threshold result,
  proposed result, and pointwise error plot.}}}
  \caption{Representative structural and geometric fitting results.}
  \label{fig:qualitative}
\end{figure}

\subsection{Limitations}

The supervised count formulation requires a specified geometric tolerance and
labeled knot vectors under a consistent parameterization. Greedy removal is
deterministic but is not guaranteed to find the globally smallest knot set.
The present study focuses on single open cubic
curves and does not yet handle closed curves, repeated knots, rational weights,
surface patch topology, or continuity constraints between CAD entities. A
an ordinal count error may still select the wrong representation, while BIC
selection increases deployment cost by refitting several complete branches.

\section{Conclusion}
\label{sec:conclusion}

We presented a framework for learning compact B-spline representations from
ordered samples by combining tolerance-aware canonical supervision, ordinal
local model-order prediction, and shared count-conditioned knot localization.
The decoder produces a valid number of ordered knots without Hard-Concrete
gates or node thresholds. A learned-prior BIC comparison of complete branches
then selects a standard regularized B-spline representation.
\TODO{Summarize the final numerical findings and the strongest demonstrated
generalization result.} Future work will extend the formulation to closed and
rational curves, repeated knots, connected curve networks, and tensor-product
surfaces.


\section*{CRediT authorship contribution statement}
\TODO{Describe each author's contributions using the CRediT taxonomy. Keep the
anonymized review version free of identifying information if required by the
submission system.}

\section*{Declaration of competing interest}
The authors declare that they have no known competing financial interests or
personal relationships that could have appeared to influence the work reported
in this paper. \TODO{Confirm this statement.}

\section*{Data availability}
The source code, trained models, and generated benchmark data will be deposited
in \TODO{repository and persistent identifier} upon publication.

\section*{Declaration of generative AI and AI-assisted technologies in the
manuscript preparation process}
During the preparation of this work, the authors used \TODO{name of tool or
remove this section} to assist with \TODO{language editing or manuscript
organization}. The authors reviewed and edited all generated material and take
full responsibility for the content of the article.

\bibliographystyle{elsarticle-num}
\bibliography{references}

\end{document}
